\section{Survivor corridors and adaptive transfer direction}
\label{sec:corridor-transfer}

Report 24 refines the full transfer of Section~\ref{sec:graded-transfer}.
It removes algebraic paths that cannot contribute to a surviving map and
chooses which end of each remaining component to propagate from.  The
maintained implementation exposes this as \code{corridor} and
\code{corridor-adaptive}; neither changes the default sparse reducer.
The latter first permits $\max(256,4N_{\rm live})$ ordinary Schur-update
pairs in a stage.  Completed pivots are retained when that allowance is
exhausted, and the remaining complex is transferred exactly.  The policy
therefore responds to observed elimination work in the spirit of introsort.
Unlike introsort's fallback theorem, it supplies neither a competitive
runtime guarantee nor a bound on the number of objects generated by a scan.

\subsection{From perturbation transfer to a directed graph}

Write $d=d_0+\delta$, with $d_0$ the weight-zero scalar differential, and
let $(I,P,H)$ be its typed binary special contraction.  Here $H$ preserves
matching type and quantum shift and lowers homological degree by one;
$\delta$ raises homological degree by one and strictly raises quantum shift.
For the degree-$h$ part of
\[
 D=P\delta(1+H\delta)^{-1}I
   =\sum_{r\ge0}P\delta(H\delta)^r I,
\]
make one vertex $z_a$ for each original degree-$h$ object and one $w_b$
for each degree-$(h+1)$ object.  Include terminal vertices $s,t$ for the
survivors in those degrees.  Its edges are precisely
\[
 s\xrightarrow{I_{as}}z_a,\qquad
 z_a\xrightarrow{\delta_{ba}}w_b,\qquad
 w_b\xrightarrow{H_{ab}}z_a,\qquad
 w_b\xrightarrow{P_{tb}}t,
\]
omitting zero entries.  The $I,H,P$ edges carry typed identities, whereas
the radical edges retain their actual cobordism coefficients.  Ordering by
quantum shift, and at equal shift by phase $w,s,z,t$, makes every edge go
forward.  Thus this graph is acyclic even when the homological degrees
of the survivors are consecutive.

\begin{proposition}[Exact survivor restriction]
The matrix entry $D_{ts}$ is the sum of the ordered coefficient products
over all directed paths from $s$ to $t$.  Deleting every vertex not both
reachable from an input and able to reach an output preserves all entries
of $D$.
\end{proposition}
\begin{proof}
Expansion of $P\delta(H\delta)^r I$ gives exactly the paths with $r$
homotopy edges.  Quantum acyclicity makes the expansion finite.  A path
contributing to any matrix entry has both required reachability properties
at each vertex, so deleting the other vertices removes none of its terms.
Conversely every retained path was already present.  An edge between two
retained vertices lies on a complete path by concatenation, so retaining
all such edges is safe.
\end{proof}

Reachability is computed by Boolean union, not by coefficient XOR.  It is
an overestimate of algebraic support: equal path products can cancel and
nonzero factors can have zero composite.  In particular, neither a nonempty
corridor nor the scalar survivor multiplicities decide a homology rank.
The implementation evaluates the entire remaining differential.

The weak components of the retained graph include its terminal vertices.
They give disjoint blocks of this adjacent-degree map.  Within a block,
forward propagation accumulates each source's coefficients in topological
order; reverse propagation accumulates each target's suffix coefficients
in reverse order.  For an edge $u\xrightarrow{f}v$, the latter forms
$g\circ f$ when the current suffix is $g:v\to t$.  It does not reverse
the order of multiplication or assume that differently typed coefficients
commute.  Distributivity and associativity show that both evaluations
give the same finite path sum.

\subsection{What the direction heuristic bounds}

Let $L(u)$ be the set of sources reaching $u$ and $R(v)$ the set of
targets reachable from $v$.  For a retained component $G_j$, define
\[
 F_j=\sum_{(u,v)\in E_j}|L(u)|,\qquad
 B_j=\sum_{(u,v)\in E_j}|R(v)|.
\]
The \code{reach} policy constructs these endpoint sets as packed integers
and chooses forward when $F_j\le B_j$, reverse otherwise.

\begin{proposition}[Propagation bound]
The number of coefficient edge propagations under this policy is at most
$\sum_j\min(F_j,B_j)$.  If component $j$ has $a_j$ inputs and $b_j$
outputs, the alternative \code{ports} policy uses only Boolean reachability
flags and bounds the propagation count by
$\sum_j |E_j|\min(a_j,b_j)$.
\end{proposition}
\begin{proof}
Forward propagation from one source processes an edge at most once, and
only if that source reaches its initial vertex.  Summing gives $F_j$.
The suffix argument gives $B_j$.  The implemented choice takes their
minimum.  Replacing the two reach-set sizes by the total numbers of ports
gives the second bound, choosing the side with fewer ports.
\end{proof}

These bounds count attempted edge propagations, not cache misses or seconds.
The code bypasses typed identities and caches equal products.  It also
scans every ordered vertex of a component once per chosen root, including
vertices whose current coefficient is zero.  If $r_j$ is the number of
roots on the chosen side, that costs $O(\sum_j |V_j|r_j)$ dictionary visits.
The \code{reach} choice need not minimize this term or choose the side with
fewer ports.  Nor does its bound include scalar setup and endpoint masks.

\subsection{Sparse scalar components and complete costs}

The scalar graph has only equal-matching, equal-shift identity arrows.
The \code{components} engine finds its weak components and contracts each
separately.  Isolated vertices are survivor identities; a component
consisting of one scalar arrow has only its reverse homotopy entry.
Other components use the existing binary contraction on local indices.
The result stores sorted tuples of nonzero indices for $I,P,H$, including
empty tuples at dead slots, instead of integers indexed by all old objects.

This changes representation without changing the established scalar basis.
Within a matching/shift/degree group, elimination sees source and target
indices in their original order.  Operations in disjoint scalar components
cannot interact.  Each selected kernel representative has a distinct largest
original source index, its dependent-column anchor.  Sorting group keys and
restoring anchor order therefore reproduces the original representative
order.  The same permutation of projection columns, and unchanged local
homotopy coefficients, recover all three original maps entrywise.

Let $N$ count allocated slots, including dead slots; let $E_{\rm in}$ count
input differential entries, and $M_{\rm sc}$ count nonzero scalar-map
indices.  If scalar components have sizes $s_i$, a conservative binary
setup bound is, up to polynomial factors in index and label bit lengths,
\[
 A_{\rm sc}=O\left(N+E_{\rm in}+M_{\rm sc}
                  +\sum_{i:s_i>2}s_i^3\right).
\]
The scan of $E_{\rm in}$ is necessary to discover scalar entries.
Across adjacent degrees, graph sizes satisfy
$V=O(N+S)$ and $E=O(E_\delta+M_{\rm sc})$, where $S\le N$ is the number
of survivors.  Graph construction, sorting, Boolean reachability, and
union by size with path compression cost near-linear work in this explicit
representation, apart from index and label arithmetic.

For the \code{ports} policy the remaining bound has the form
\[
 T\le A_{\rm sc}+A_{\rm graph}
       +\Lambda\sum_j |V_j|r_j
       +\Lambda(1+C(k))\sum_j |E_j|r_j,
 \qquad r_j=\min(a_j,b_j),
\]
where $\Lambda$ charges index/dictionary arithmetic and $C(k)$ bounds
coefficient work at $2k$ boundary points.  A conservative coefficient
envelope is $\operatorname{poly}(k)2^{O(k)}$.  The \code{reach} version
replaces the edge term by $\sum_j\min(F_j,B_j)$, but uses its actually
chosen $r_j$ in the vertex term and additionally pays up to
$O((V+E)S)$ bit work for endpoint masks and population counts.
Python hash-table timings are empirical; ordered dictionaries give a
deterministic accounting with polynomial index overhead.

Sparse indices matter even for a singleton column: an integer $2^j$ uses
$j+1$ bits of payload.  Many such global columns can have quadratic total
bit length while their sparse representation has only linearly many
indices.  Conversely, a dense connected scalar block can still require
large local binary matrices.  This optimization does not remove all fill.

\subsection{A strict separation, and its scope}

Fix the three-arc dot algebra $\F[x,y,z]/(x^2,y^2,z^2)$.
Take $R$ degree-zero survivors $s_j$ at shift zero and one degree-one
survivor $t$ at shift six.  Add identity pairs $L_0\to B_0$ at shift two
and $L_i\to B_i$ at shift four for $1\le i\le M$, with $L$ in degree
zero and $B$ in degree one.  The only radical maps are
\[
 s_j\to B_0:x,\qquad L_0\to B_i:y,\qquad L_i\to t:z.
\]
This is a homogeneous two-term complex, hence $d^2=0$.  For odd $M$ its
transferred entry is $D_{ts_j}=Mxyz=xyz\ne0$.  Forward propagation visits
$R(1+2M)$ radical edges, while reverse propagation visits $R+2M$.
The latter accumulates the common suffix once before distributing it to
the $R$ sources.  For even $M$ the output vanishes by parity despite
unchanged structural reachability, as a regression test checks.

There are $R+2M+3$ original objects, $R+1$ isolated scalar survivors, and
$M+1$ identity pairs.  Sparse scalar setup has no nontrivial binary blocks.
Using \code{components} and \code{ports}, the single-target graph is
evaluated once in reverse; all coefficients lie in a fixed eight-dimensional
algebra.  Including input reading, scalar setup, sorting and output writing,
the stage takes $O((R+M)\log(R+M+2))$ index-word operations and
$O(R+M)$ words.  The previous forward evaluation alone takes $\Omega(RM)$
propagations.  At $R=M=255$, the radical counts are $130305$ and $765$;
the contraction uses $768$ sparse indices versus $164735$ packed payload
bits.  These are representation counts, not heap measurements.

This is a separation between complete transfer implementations on an
explicit algebraic family.  Its realization by knot prefixes is unproved,
and minimum-fill sparse cancellation is already efficient on the family.
The default public corridor policy still uses \code{reach} and the global
scalar engine; the near-linear theorem refers specifically to the direct
\code{ports}/\code{components} configuration.

\subsection{Integration, validation, and measured limits}

The new reducers support the standard scanner, minimum-fill pivots,
bit coefficients, the existing component coefficient engines, and tail
contraction.  Races, other scanner backends, and window truncation are
rejected before early certificates.  Window integration needs its own
quantum-shift and truncation argument.  Global resource exhaustion raises
the ordinary scan limit in the raw API and yields \code{UNKNOWN} in
recognition.  Full replacement arrays, reverse indices, and an empty cache
are allocated before a final deadline check and installation.  Tests inject
both allocation failure and expiry after transfer, then resume successfully
on the preserved complex.  Completed earlier sparse pivots remain valid.

The maintained independent audit compares seven configurations entrywise
against the frozen forward reference on $555$ prefixes from $80$ diagrams:
$3885$ comparisons, including partially cancelled inputs.  All eight special
contraction identities are checked at each prefix.  Separate tests compare
graded homology, continued scans and independent set coefficients.  The
full maintained suite passes $458$ tests.  The
stronger graded support-only certificate uses the condition
\[
 \Delta>0,\quad k-c(a,b)\le\Delta\le k+c(a,b),\quad
 \Delta\equiv k-c(a,b)\pmod2
\]
for a possible radical map.  Its audit accepts $196$ of $639$ stages in
$72$ diagrams, exactly the same stages as the older degree-gap shortcut.
The experimental support reducer is retained for comparison, not substituted
for an established production policy.

The delivered package lacks its advertised full SHA-256 manifest.  The
maintained reproduction wrapper verifies the separate $20$-file integration
manifest and runs the archive in a temporary copy: $268$ scanner tests,
the transfer and support audits, and the independent surface-cover suite
all pass.  This is a narrower provenance claim than verifying every file.
The surface-cover theorem has a separate presentation contract and is not
used to make a knot verdict in this integration.

\begin{center}
\begin{tabular}{@{}lrrrrr@{}}
\toprule
Raw scan & Std/corridor & Std/adapt. & Std/sparse & A/A & Switches \\
\midrule
two\_strand\_201 & 0.105 & 0.808 & 0.084 & 1.019 & 0 \\
torus\_3\_10 & 0.327 & 0.932 & 0.305 & 0.987 & 0 \\
alternating\_3\_4 & 0.348 & 0.938 & 0.302 & 1.017 & 0 \\
morton & 0.422 & 0.927 & 0.411 & 1.013 & 0 \\
conway & 0.341 & 0.923 & 0.309 & 0.991 & 0 \\
hard\_unknot\_8 & 0.500 & 0.918 & 0.490 & 0.986 & 0 \\
\bottomrule
\end{tabular}
\end{center}
Sparse denotes the component-scalar/Boolean-port configuration. Ratios above one favor the denominator.
\begin{center}
\begin{tabular}{@{}lrrrr@{}}
\toprule
Constructed stage & Fwd/corridor & Fwd/sparse & Std/sparse & A/A \\
\midrule
source\_heavy & 2.390 & 2.214 & 0.340 & 0.988 \\
sink\_heavy & 0.461 & 0.426 & 0.336 & 1.016 \\
balanced & 0.623 & 0.592 & 0.299 & 1.014 \\
dead\_branches & 6.234 & 5.457 & 1.188 & 0.993 \\
shared\_suffix\_15 & 1.303 & 1.558 & 0.204 & 1.017 \\
shared\_suffix\_63 & 4.194 & 4.859 & 0.172 & 0.981 \\
shared\_suffix\_255 & 15.661 & 18.141 & 0.154 & 1.041 \\
dead\_leaves\_255 & 0.677 & 0.860 & 0.187 & 0.993 \\
\bottomrule
\end{tabular}
\end{center}
Fwd denotes the maintained quantum-ordered forward transfer. These timings include scalar and graph setup.


All six actual scans favor ordinary sparse cancellation: eager corridor
transfer is approximately $2.0$--$9.5$ times slower.  The adaptive policy
switches zero times, with standard/adaptive ratios $0.808$--$0.938$;
its quantum bookkeeping therefore buys no transfer benefit in these cases.
The sparse-scalar/port refinement also loses on every actual scan.
Actual standard/control ratios span $0.986$--$1.019$.
On the constructed size-$255$ shared-suffix case, the sparse refinement
is $18.14$ times faster than the maintained forward transfer, but ordinary
sparse cancellation is still approximately $6.49$ times faster than that
refinement.  The constructed dense dead-branch case gives a different
result: eager corridor and adaptive corridor beat ordinary cancellation
by factors $1.406$ and $1.380$, respectively.  Other constructed cases
include regressions.  These observations support retaining optional
policies and publishing their negative results, rather than changing the
default or claiming a uniform timing improvement.

The measurements use fresh states, seven shuffled paired warm rounds,
and identical standard/control arms.  Actual raw scans include construction
of the scanner, quantum tracking and all reduction work; diagram construction
and shared crossing-order choice are outside timing.  Synthetic timings
include the complete reduction stage, with fixture construction outside
timing.  Early recognition filters are bypassed, so these are not
end-to-end recognition speedups.  Raw samples, repetition counts, exact
checks and source hashes are retained with the result file.

Finally, a scan with every allocated $N_i\le\exp(O(\log^2(n+2)))$ and
every $k_i=O(\log^2(n+2))$, suitably bounded labels and an affordable order,
would meet the target by the displayed bounds and ordinary extension costs.
No such uniform bounds have been proved for arbitrary diagrams.  Corridor
size also depends on the scalar basis; it is not a knot invariant.  Faster
evaluation of an explicit stage leaves the central multiplicity and
representation problems unresolved.

The weaker general sub-exponential target is also unproved for this
implementation.  Here sub-exponential means $2^{o(n)}$ in the crossing
count, not merely a speedup over the previous program.  Under the same
polynomial label-size assumptions, the stage accounting would suffice
if both $\max_i k_i=o(n)$ and $\max_i\log(N_i+2)=o(n)$, with an order
found in sub-exponential time.  No such simultaneous guarantee has been
established.  Adaptive switching among the current representations does
not supply it.  This implementation status must also be distinguished
from Lackenby's announced quasi-polynomial result: the July 2026 hierarchy
paper explicitly leaves some step runtimes unspecified in its
Section~9~\cite{lackenby2026}.
